% Lösungen Einheit 1
\einheitenkopf{Einheit 1 von 5 · Monotonie und erste Ableitung}
\erg{8}{a) steigt \quad b) fällt \quad c) steigt \quad d) fällt}
\erg{9}{a) $f'(x) = x^2 - 9 = 0 \Leftrightarrow x = \pm 3$; steigt auf $]{-\infty};\,-3]$ und $[3;\,\infty[$, fällt auf $[-3;\,3]$ \quad b) $f'(x) = -3(x + 1)(x - 3)$, $f'(0) = 9 > 0$: fällt auf $]{-\infty};\,-1]$ und $[3;\,\infty[$, steigt auf $[-1;\,3]$ \quad c) $f'(x) = 4x(x^2 - 4)$; Extremstellen $-2$, $0$, $2$; fällt auf $]{-\infty};\,-2]$ und $[0;\,2]$, steigt auf $[-2;\,0]$ und $[2;\,\infty[$ \quad d) steigt auf $]{-\infty};\,-1]$ und $[3;\,\infty[$, fällt auf $[-1;\,3]$ \quad e) $f'(x) = (x - 1)\,\mathrm{e}^{x}$; fällt auf $]{-\infty};\,1]$, steigt auf $[1;\,\infty[$ \quad f) $d(t) = -0{,}1t^3 + 1{,}5t^2 - 2{,}7t$; $d'(t) = -0{,}3(t - 1)(t - 9) = 0 \Leftrightarrow t = 1$ oder $t = 9$; $d$ fällt auf $[0;\,1]$ und $[9;\,10]$, steigt auf $[1;\,9]$; $d(0) = 0$, $d(1) = -1{,}3$, $d(9) = 24{,}3$, $d(10) = 23$; Wertebereich $[-1{,}3;\,24{,}3]$ (in m³)}
\erg{10}{a) $f'(x) = 3x^2 + 5 \ge 5 > 0$: streng monoton steigend \quad b) $f'(x) = 1 + \mathrm{e}^{x} > 1$: streng monoton steigend \quad c) $f'(x) = -3x^2 - 2 \le -2 < 0$: streng monoton fallend \quad d) $f'(x) \ge 5$ wird nie null; ohne waagerechte Tangente kein Hoch- oder Tiefpunkt \quad e) Produktregel: $f'(x) = 2x\,\mathrm{e}^{x} + (x^2 + 1)\,\mathrm{e}^{x} = (x^2 + 2x + 1)\,\mathrm{e}^{x} = (x + 1)^2\,\mathrm{e}^{x}$; $(x + 1)^2 \ge 0$ und $\mathrm{e}^{x} > 0$, also $f'(x) \ge 0$: der Graph fällt nirgends}
\erg{11}{a) Mia untersucht das Vorzeichen von $f''$ statt von $f'$. Richtig: $f'(x) = 3x(x - 4) = 0 \Leftrightarrow x = 0$ oder $x = 4$; $f'(-1) = 15 > 0$, $f'(1) = -9 < 0$, $f'(5) = 15 > 0$: steigt auf $]{-\infty};\,0]$ und $[4;\,\infty[$, fällt auf $[0;\,4]$ \quad b) $g'(x) = 3x(x + 2)$: steigt auf $]{-\infty};\,-2]$ und $[0;\,\infty[$, fällt auf $[-2;\,0]$}
\erg{12}{a) $4x^2 \ge 0$, also $f'(x) \ge 7 > 0$; ein Hoch- oder Tiefpunkt bräuchte $f'(x) = 0$. \quad b) Tim irrt: $f'(x) = 5x^4 > 0$ für alle $x \ne 0$; an der einzelnen Stelle $0$ ist die Tangente waagerecht, der Graph fällt aber nirgends – $f$ ist streng monoton steigend.}
\erg{13}{a) $A'(t) = -0{,}06t\,(t - 10) \ge 0$ auf $[0;\,10]$: nach Modell A steigt die Temperatur von 6 bis 16 Uhr; $B'(t) = -0{,}2t + 1{,}6 \ge 0$ auf $[0;\,8]$: nach Modell B von 6 bis 14 Uhr \quad b) Die Aussage stimmt: 10 Stunden gegen 8 Stunden.}
